\chapter{杂环化学}\label{ch:b3:heterocycles}

一杯咖啡中含有咖啡因，咖啡因由两个稠合环构成，这两个环共含
四个氮原子；咖啡香气则来自焙烧过程中生成的数十种呋喃、吡咯、吡嗪
和噻吩。药房货架上的大多数药物至少含有一个
环，该环中含氮、氧或硫原子，DNA 的碱基、血红素、叶绿素及
许多维生素也是如此。本章阐述芳香环中的杂原子如何改变环上的电子：
它可使环比苯“缺电子”或“富电子”，
由此决定其反应的位置与方式。
本章最后介绍三种经典的此类环合成方法。

\begin{recall}
本科二年级分册讲授了芳香性与休克尔规则、芳烃及
经惠兰德中间体的亲电芳香取代、活化与定向基团、胺及其碱性、
亚胺、烯胺、互变异构体与碱基；本科一年级分册讲授了 $\mathrm pK_a$、
亲核试剂、亲电试剂与离去基团。
\Cref{ch:b3:pericyclic}中给出了 [3,3]
σ 迁移重排。
\end{recall}

\begin{omfigure}
\includegraphics[width=0.5\linewidth]{images/book4/ai/b3-29-coffee.jpg}

{\small 一杯咖啡与焙烧后的咖啡豆：咖啡因（一种嘌呤）以及
焙烧过程中生成的许多香气分子都是杂环化合物。}
\end{omfigure}

\section{芳香杂环}

\begin{definition}[杂环]\label{def:b3:heterocycles:heterocycle}
称\emph{杂环}\index{杂环}为环中至少含有一个
非碳原子的环状化合物；称\emph{杂芳香化合物}
\index{杂芳香化合物}为其中的芳香化合物，例如吡啶、
吡咯、呋喃与噻吩。
\end{definition}

\begin{proposition}[六个 $\pi$ 电子]\label{prop:b3:heterocycles:sextet}
吡啶、吡咯、呋喃、噻吩与咪唑都在一个平面、环状、共轭的环中
具有六个 $\pi$ 电子，按休克尔规则均为芳香性的。
\end{proposition}

\begin{proof}
吡啶：五个碳原子和氮原子各提供一个 p 电子；氮原子的
孤电子对处于环平面内的 $sp^2$ 轨道中，不在 $\pi$ 体系内：$5 + 1 = 6$。
吡咯：四个碳原子各提供一个电子，与平面内 H 相连的氮原子
将其孤电子对放入 p 轨道：$4 + 2 = 6$。呋喃与噻吩与吡咯类似，杂原子
提供两对孤电子对中的一对（另一对位于平面内）。咪唑：
N--H 氮原子提供两个电子，类吡啶型氮原子提供一个，三个碳原子
各提供一个，共 6 个。6 符合 $4n + 2$，此时 $n = 1$。
\end{proof}

\begin{omfigure}
\begin{tikzpicture}[font=\small, yscale=0.55]
  % pyridine: six p orbitals perpendicular to the ring (drawn first), N at 0 deg
  \foreach \a in {0,60,...,300} { \pic at (\a:1.1) {omp={0.62}{90}}; }
  \draw[thick] (0:1.1) -- (60:1.1) -- (120:1.1) -- (180:1.1) -- (240:1.1) -- (300:1.1) -- cycle;
  \node[fill=white, inner sep=0.5pt] at (0:1.1) {N};
  \pic at ($(0:1.1)+(0.2,0)$) {omlobe={0.8}{0}};
  \node[font=\scriptsize, align=center, anchor=north] at (0,-2.4) {吡啶：6 个原子提供 6 个 p 电子；\\ 孤电子对位于平面内（具碱性）};
  \begin{scope}[xshift=5.4cm]
    \foreach \a in {0,72,144,216,288} { \pic at (\a:1.0) {omp={0.62}{90}}; }
    \draw[thick] (0:1.0) -- (72:1.0) -- (144:1.0) -- (216:1.0) -- (288:1.0) -- cycle;
    \node[fill=white, inner sep=0.5pt] at (0:1.0) {N};
    \draw ($(0:1.0)+(0.18,0)$) -- (0:1.7) node[right, inner sep=0.5pt] {H};
    \node[font=\scriptsize, align=center, anchor=north] at (0,-2.4) {吡咯：N 的孤电子对是 \\ 6 个 $\pi$ 电子之一（无碱性）};
  \end{scope}
\end{tikzpicture}

{\small 吡啶与吡咯的 $\pi$ 体系（斜视）：每个环原子各有一个
与环面垂直的 p 轨道。在吡啶中，氮原子上位于平面内的
孤电子对（瓣）沿环平面向外伸出，可以结合质子；在吡咯中，
氮原子带有平面内的氢，其孤电子对贡献给
芳香六电子体系。}
\end{omfigure}

\begin{definition}[多 $\pi$ 杂环与缺 $\pi$ 杂环]\label{def:b3:heterocycles:pi-character}
称\emph{多 $\pi$ 杂环}\index{多 $\pi$ 杂环}，如吡咯、呋喃或噻吩，
五个原子上共有六个 $\pi$ 电子，其碳原子上的
$\pi$ 电子密度比苯高；称\emph{缺 $\pi$ 杂环}
\index{缺 $\pi$ 杂环}，例如吡啶，具有一个
电负性较大的环原子，使 $\pi$ 电子密度自碳原子处被拉走。
\end{definition}

\begin{proposition}[碱性]\label{prop:b3:heterocycles:basicity}
吡啶为中等强度的碱，吡咯几乎不具碱性；吡咯在
强酸中被质子化时，质子结合在碳原子而非氮原子上。
\end{proposition}

\begin{proof}[论证]
% ledger: pka:pyridinium, pka4:pyrrolium, pka4:piperidinium, pka4:imidazolium, pka4:DMAPH
吡啶的孤电子对位于 $\pi$ 体系之外，质子化它不会损失
芳香性，吡啶鎓的 $\mathrm pK_a = 5.2$，其碱性弱于哌啶
（吡啶鎓的氮原子为 $sp^2$ 杂化，s 成分更多，对其孤电子对的束缚
更紧：哌啶鎓为 11.1）。若将吡咯的氮原子质子化，则会使其孤电子对
脱离六电子体系，破坏芳香性；只有很强的酸
才能使吡咯发生质子化，质子结合在 C2 位上，此时正离子
保留一定离域作用（$\mathrm pK_a$ 约为 $-4$）。咪唑两类氮原子兼有，它在类吡啶型氮原子上
被质子化，其阳离子对称，通过共振得以稳定（$\mathrm pK_a$
7.0），这也是组氨酸在酶中能兼作酸和碱的原因。
4-（二甲氨基）吡啶的碱性更强（$\mathrm pK_a$ 9.6）：氨基将其
孤电子对推入环中，并传递到环氮原子上。
\end{proof}

% ledger: dip:pyridine, dip4:pyrrole, dip4:furan, dip4:thiophene
偶极矩同样说明了这一点：吡啶（\qty{2.19}{D}）的负端位于
氮原子上；而在吡咯（\qty{1.84}{D}）中，孤电子对向环的给予
使负端位于环上、正端位于氮原子上。在呋喃
（\qty{0.66}{D}）与噻吩（\qty{0.55}{D}）中，$\sigma$ 吸电子与 $\pi$ 给电子两种效应，几乎
相互抵消。

\begin{omfigure}
% ledger: h1:pyridine.a, h1:pyridine.g, h1:pyridine.b, h1:pyrrole.a, h1:pyrrole.b, h1:pyrrole.NH, h1:furan.a, h1:furan.b
\begin{tikzpicture}
\begin{axis}[width=0.86\linewidth, height=6.2cm, x dir=reverse, xmin=5.8, xmax=9.0, ymin=-0.1, ymax=3.6,
  axis x line=bottom, axis y line=none, xlabel={化学位移 $\delta$ / ppm}, tick label style={font=\scriptsize}, label style={font=\small}]
\addplot[omDef, thick] table[x=ppm, y expr=\thisrow{y}+2.4] {figdata/out/bachelor-3-heterocycle-nmr-pyridine.dat};
\addplot[omThm, thick] table[x=ppm, y expr=\thisrow{y}+1.2] {figdata/out/bachelor-3-heterocycle-nmr-pyrrole.dat};
\addplot[omMeth, thick] table[x=ppm, y=y] {figdata/out/bachelor-3-heterocycle-nmr-furan.dat};
\node[font=\scriptsize, omDef, anchor=west] at (axis cs:7.1,3.0) {吡啶：H2/6、H4、H3/5（2\,:\,1\,:\,2）};
\node[font=\scriptsize, omThm, anchor=west] at (axis cs:7.7,1.8) {吡咯：NH、H2/5、H3/4};
\node[font=\scriptsize, omMeth, anchor=west] at (axis cs:7.2,0.6) {呋喃：H2/5、H3/4};
\end{axis}
\end{tikzpicture}

{\small 吡啶、吡咯与呋喃在氘代氯仿中的质子核磁共振谱，
由实测化学位移绘制为单线（略去耦合），
各峰面积等于质子数。缺 $\pi$ 的吡啶，其 H2/6 与 H4
信号明显位于苯的低场；多 $\pi$ 的吡咯与呋喃的大部分
其信号位于苯的高场一侧；与杂原子相邻的质子始终
去屏蔽程度最大。}
\end{omfigure}

\section{吡啶}

吡啶难以发生亲电取代：亲电试剂所遇到的环
电子比苯更少，且在通常存在的酸中，生成
带正电荷的吡啶正离子，其正电荷排斥亲电试剂。硝化需要非常苛刻的
条件，且以较低产率得到 3-异构体。

\begin{proposition}[C3 位上的亲电取代]\label{prop:b3:heterocycles:pyridine-c3}
吡啶的亲电进攻发生在 C3 位：进攻 C2 或 C4 位得到的
惠兰德中间体，其一个共振结构使正电荷位于氮原子上，
而该氮原子周围只有六个电子。
\end{proposition}

\begin{proof}
进攻 C2（或 C4）所得正离子中，正电荷分布于 C3、C5
与 N1（或 C3、C5 与 N1）上；其一个共振结构含缺电子的
六电子氮原子，这对电负性原子极为不利。进攻 C3 时
正电荷仅分布于 C2、C4 与 C6。三种正离子的稳定性
均低于苯的，但 C3 进攻的不利程度最小。
\end{proof}

\begin{definition}[吡啶 N-氧化物]\label{def:b3:heterocycles:n-oxide}
由过氧酸氧化吡啶所得的\emph{吡啶 N-氧化物}\index{吡啶 N-氧化物}，
带有 $\mathrm{N^+{-}O^-}$ 基团，其氧原子向 C2
和 C4 给电子；它在 C4 位发生亲电取代，随后
脱去氧原子。
\end{definition}

\begin{definition}[亲核芳香取代]\label{def:b3:heterocycles:snar}
称\emph{亲核芳香取代}\index{亲核芳香取代}
（$\mathrm S_\mathrm NAr$）为以亲核试剂取代缺电子芳环上离去基团的
反应，历程为先加成、后消除。其阴离子形式的
加成中间体即\emph{迈森海默配合物}\index{迈森海默配合物}。
\end{definition}

\begin{proposition}[在 C2 与 C4 位的 $\mathrm S_\mathrm NAr$]\label{prop:b3:heterocycles:snar-positions}
卤代吡啶在 C2 与 C4 位上容易发生 $\mathrm S_\mathrm NAr$，在 C3 位则
进行得十分缓慢，因为在 C2 或 C4 位加成时，负电荷位于氮原子上。
\end{proposition}

\begin{proof}
在 C2 位加成使 C2 成为 $sp^3$ 杂化，余下的四个 $\pi$ 电子与
进入的一对电子分布于 N1、C3 与 C5（若为 C4，则为 C3、C5 与 N1）；
其一个共振结构使负电荷位于电负性大的氮原子上。在 C3 位加成时，
负电荷仅分布于 C2、C4 与 C6，均为碳原子。
\end{proof}

\begin{omfigure}
\begin{tikzpicture}[font=\small, >=Stealth]
  \draw[thick] (-0.000,-0.620) -- (0.537,-0.310) -- (0.537,0.310) -- (0.000,0.620) -- (-0.537,0.310) -- (-0.537,-0.310) -- cycle;
  \draw[thick] (0.082,-0.424) -- (0.326,-0.283);
  \draw[thick] (0.326,0.283) -- (0.082,0.424);
  \draw[thick] (-0.408,0.141) -- (-0.408,-0.141);
  \node[fill=white, inner sep=0.4pt, font=\scriptsize] at (-0.000,-0.620) {N};
  \draw (0.537,-0.310) -- (0.927,-0.535) node[right, inner sep=0.5pt, font=\scriptsize] {Cl};
  \draw[thick] (3.400,-0.620) -- (3.937,-0.310) -- (3.937,0.310) -- (3.400,0.620) -- (2.863,0.310) -- (2.863,-0.310) -- cycle;
  \draw[thick] (3.726,0.283) -- (3.482,0.424);
  \draw[thick] (2.992,0.141) -- (2.992,-0.141);
  \node[fill=white, inner sep=0.4pt, font=\scriptsize] at (3.400,-0.620) {N};
  \node[font=\scriptsize, omThm] at (3.350,-0.920) {$\ominus$};
  \draw (3.937,-0.310) -- (4.386,-0.343) node[right, inner sep=0.5pt, font=\scriptsize] {Nu};
  \draw (3.937,-0.310) -- (4.190,-0.682) node[right, inner sep=0.5pt, font=\scriptsize] {Cl};
  \draw[thick] (6.800,-0.620) -- (7.337,-0.310) -- (7.337,0.310) -- (6.800,0.620) -- (6.263,0.310) -- (6.263,-0.310) -- cycle;
  \draw[thick] (6.882,-0.424) -- (7.126,-0.283);
  \draw[thick] (7.126,0.283) -- (6.882,0.424);
  \draw[thick] (6.392,0.141) -- (6.392,-0.141);
  \node[fill=white, inner sep=0.4pt, font=\scriptsize] at (6.800,-0.620) {N};
  \draw (7.337,-0.310) -- (7.727,-0.535) node[right, inner sep=0.5pt, font=\scriptsize] {Nu};
  \draw[->, thick] (1.25,0) -- (2.35,0) node[midway, above, font=\scriptsize] {$+\,\mathrm{Nu^-}$};
  \draw[->, thick] (4.65,0) -- (5.75,0) node[midway, above, font=\scriptsize] {$-\,\ce{Cl-}$};
  \node[font=\scriptsize] at (3.4,-1.2) {迈森海默配合物};
\end{tikzpicture}

{\small 2-氯吡啶的亲核取代：亲核试剂加成至
C2 位，生成阴离子中间体，其负电荷主要位于
环氮原子上（图中示出一种共振结构）；随后氯离子离去，
芳环得以恢复。}
\end{omfigure}

\begin{proposition}[加成–消除]\label{prop:b3:heterocycles:snar-mechanism}
对于加成步骤为决速步的 $\mathrm S_\mathrm NAr$ 反应，$v = k[\mathrm{ArX}][\mathrm{Nu^-}]$，
且氟是最好的卤素离去基团。
\end{proposition}

\begin{proof}[论证]
反应速率即双分子加成的速率。在该步骤中，碳–卤键
并未断裂；关键在于卤素通过吸电子降低阴离子过渡态
能量的程度，氟的电负性最大，
其作用最强：F $>$ Cl $\approx$ Br $>$ I，与 $\mathrm S_\mathrm N2$ 的顺序相反。
\end{proof}

\begin{definition}[齐齐巴宾反应]\label{def:b3:heterocycles:chichibabin}
\emph{齐齐巴宾反应}\index{齐齐巴宾反应}是用氨基钠对吡啶的
C2 位进行的氨基化：酰胺负离子加成到 C2 上，并失去一个
氢负离子，该氢负离子与产物反应后以氢气的形式放出。
\end{definition}

吡啶和 DMAP 也用作碱和亲核催化剂：在
与酸酐的酰化反应中，DMAP 先进攻酸酐，生成一个
酰基吡啶正离子，它把酰基转移到醇上的速度远快于
酸酐本身。吡啶、联吡啶和菲啰啉是配位化学中最常见的
配体之一。

\section{五元环}

\begin{proposition}[吡咯在 C2 位的亲电取代]\label{prop:b3:heterocycles:pyrrole-c2}
吡咯、呋喃和噻吩主要在 C2 位发生亲电取代，
原因在于进攻 C2 位所生成的阳离子有三种共振结构，
而进攻 C3 位的只有两种。
\end{proposition}

\begin{proof}
进攻 C2 位后，正电荷可位于 C3、C5 或氮原子上，
此时正电荷通过氮的孤电子对离域，且每个原子都保持八隅体（第三种
结构）。进攻 C3 位后，正电荷只能位于 C2 或氮上，电荷
的离域程度较低（见图）。
\end{proof}

\begin{omfigure}
\begin{tikzpicture}[font=\small]
  \draw[thick] (0.000,0.620) -- (0.590,0.192) -- (0.364,-0.502) -- (-0.364,-0.502) -- (-0.590,0.192) -- cycle;
  \draw[thick] (-0.303,-0.269) -- (-0.403,0.039);
  \node[fill=white, inner sep=0.4pt, font=\scriptsize] at (0.000,0.620) {N};
  \draw (0.000,0.750) -- (0.000,1.020) node[above, inner sep=0.5pt, font=\scriptsize] {H};
  \draw (0.590,0.192) -- (0.893,0.482) node[right, inner sep=0.5pt, font=\scriptsize] {H};
  \draw (0.590,0.192) -- (1.006,0.135) node[right, inner sep=0.5pt, font=\scriptsize] {E};
  \node[font=\scriptsize, omThm] at (0.164,-0.226) {$\oplus$};
  \draw[thick] (2.600,0.620) -- (3.190,0.192) -- (2.964,-0.502) -- (2.236,-0.502) -- (2.010,0.192) -- cycle;
  \draw[thick] (2.762,-0.371) -- (2.438,-0.371);
  \node[fill=white, inner sep=0.4pt, font=\scriptsize] at (2.600,0.620) {N};
  \draw (2.600,0.750) -- (2.600,1.020) node[above, inner sep=0.5pt, font=\scriptsize] {H};
  \draw (3.190,0.192) -- (3.493,0.482) node[right, inner sep=0.5pt, font=\scriptsize] {H};
  \draw (3.190,0.192) -- (3.606,0.135) node[right, inner sep=0.5pt, font=\scriptsize] {E};
  \node[font=\scriptsize, omThm] at (2.335,0.086) {$\oplus$};
  \draw[thick] (5.200,0.620) -- (5.790,0.192) -- (5.564,-0.502) -- (4.836,-0.502) -- (4.610,0.192) -- cycle;
  \draw[thick] (5.362,-0.371) -- (5.038,-0.371);
  \draw[thick] (5.113,0.395) -- (4.851,0.205);
  \node[fill=white, inner sep=0.4pt, font=\scriptsize] at (5.200,0.620) {N};
  \draw (5.200,0.750) -- (5.200,1.020) node[above, inner sep=0.5pt, font=\scriptsize] {H};
  \draw (5.790,0.192) -- (6.093,0.482) node[right, inner sep=0.5pt, font=\scriptsize] {H};
  \draw (5.790,0.192) -- (6.206,0.135) node[right, inner sep=0.5pt, font=\scriptsize] {E};
  \node[font=\scriptsize, omThm] at (5.480,0.700) {$\oplus$};
  \draw[thick] (8.400,0.620) -- (8.990,0.192) -- (8.764,-0.502) -- (8.036,-0.502) -- (7.810,0.192) -- cycle;
  \draw[thick] (8.097,-0.269) -- (7.997,0.039);
  \node[fill=white, inner sep=0.4pt, font=\scriptsize] at (8.400,0.620) {N};
  \draw (8.400,0.750) -- (8.400,1.020) node[above, inner sep=0.5pt, font=\scriptsize] {H};
  \draw (8.764,-0.502) -- (9.135,-0.700) node[right, inner sep=0.5pt, font=\scriptsize] {H};
  \draw (8.764,-0.502) -- (8.839,-0.915) node[right, inner sep=0.5pt, font=\scriptsize] {E};
  \node[font=\scriptsize, omThm] at (8.665,0.086) {$\oplus$};
  \draw[thick] (11.000,0.620) -- (11.590,0.192) -- (11.364,-0.502) -- (10.636,-0.502) -- (10.410,0.192) -- cycle;
  \draw[thick] (10.697,-0.269) -- (10.597,0.039);
  \draw[thick] (11.087,0.395) -- (11.349,0.205);
  \node[fill=white, inner sep=0.4pt, font=\scriptsize] at (11.000,0.620) {N};
  \draw (11.000,0.750) -- (11.000,1.020) node[above, inner sep=0.5pt, font=\scriptsize] {H};
  \draw (11.364,-0.502) -- (11.735,-0.700) node[right, inner sep=0.5pt, font=\scriptsize] {H};
  \draw (11.364,-0.502) -- (11.439,-0.915) node[right, inner sep=0.5pt, font=\scriptsize] {E};
  \node[font=\scriptsize, omThm] at (11.280,0.700) {$\oplus$};
  \node[font=\scriptsize] at (2.6,-1.35) {进攻 C2 位：三种结构};
  \node[font=\scriptsize] at (9.7,-1.35) {进攻 C3 位：两种结构};
  \draw[black!40] (6.8,-1.2) -- (6.8,1.2);
\end{tikzpicture}

{\small 吡咯的惠兰德中间体（E：亲电试剂；$\oplus$：正电荷）。
进攻 C2 位生成的阳离子，其电荷位于 C3、C5 或氮上，且所有
八隅体均完整；进攻 C3 位时，电荷只能位于 C2 或氮上。}
\end{omfigure}

吡咯的反应活性与苯酚或苯胺相当：在温和条件下即可溴代、硝化
和酰化，在强酸中会聚合，并且
在 C2 位被由二甲基甲酰胺和三氯氧磷
制得的维尔斯迈尔试剂甲酰化。反应活性顺序为吡咯 $>$ 呋喃 $>$ 噻吩 $>$ 苯，
取决于杂原子提供孤电子对参与共享的能力。呋喃
芳香性最弱，在第尔斯–阿尔德反应中也可作双烯体参与反应。强
碱如丁基锂，可使呋喃和噻吩在紧邻
杂原子的 C2 位脱质子，生成有机锂试剂。

\begin{method}[预测杂环上的进攻位点]\label{met:b3:heterocycles:site}
\begin{enumerate}
  \item 对杂环进行分类：多 $\pi$ 杂环（类吡咯型杂原子）或
        缺 $\pi$ 杂环（类吡啶型氮）。
  \item 亲电试剂：作用于多 $\pi$ 杂环的 C2 位（吲哚取 C3 位，因其
        C2 位进攻会破坏苯环）；吡啶则在 C3 位，或经由
        吡啶 N-氧化物在 C4 位反应。
  \item 亲核试剂：作用于缺 $\pi$ 杂环的 C2 或 C4 位，尤其是该位置带有
        离去基团时。
  \item 强碱：作用于紧邻杂原子的 C--H。
\end{enumerate}
\end{method}

\section{稠杂环与生物杂环}

吲哚是苯与吡咯稠合而成的化合物，与亲电试剂在 C3 位反应：此时
中间体保持苯环完整，电荷位于氮上。
喹啉是苯与吡啶稠合而成，其苯环上发生亲电取代，
亲核取代则发生在 C2 和 C4 位。碱基为
嘧啶（胞嘧啶、胸腺嘧啶、尿嘧啶）和嘌呤（腺嘌呤、鸟嘌呤），
嘌呤是嘧啶与咪唑稠合而成。它们的氧代形式和氨基形式（内酰胺和胺类
互变异构体）比羟基形式和亚胺形式更占优势，
从而决定了每个碱基上氢键给体与受体的分布，进而决定碱基配对。

\begin{omfigure}
\begin{tikzpicture}[font=\small]
  \node[draw, rounded corners, minimum width=1.8cm, minimum height=2.8cm, fill=omDef!8] at (0,0) {};
  \node[font=\small\bfseries] at (0,1.75) {鸟嘌呤};
  \node[anchor=east] (g1) at (0.85,0.85) {C6=O};
  \node[anchor=east] (g2) at (0.85,0) {N1--H};
  \node[anchor=east] (g3) at (0.85,-0.85) {C2--N--H};
  \node[draw, rounded corners, minimum width=1.8cm, minimum height=2.8cm, fill=omThm!8] at (3.2,0) {};
  \node[font=\small\bfseries] at (3.2,1.75) {胞嘧啶};
  \node[anchor=west] (c1) at (2.35,0.85) {H--N--C4};
  \node[anchor=west] (c2) at (2.35,0) {N3};
  \node[anchor=west] (c3) at (2.35,-0.85) {O=C2};
  \draw[omMeth, very thick, dotted] (g1.east) -- (c1.west);
  \draw[omMeth, very thick, dotted] (g2.east) -- (c2.west);
  \draw[omMeth, very thick, dotted] (g3.east) -- (c3.west);
\end{tikzpicture}

{\small 鸟嘌呤–胞嘧啶碱基对示意图：三个氢键（绿色）
分别位于鸟嘌呤的羰基氧与胞嘧啶的一个氨基 N--H 之间，鸟嘌呤的 N1--H
与胞嘧啶的环 N3 之间，以及鸟嘌呤的一个氨基 N--H
与胞嘧啶的羰基氧之间。腺嘌呤与胸腺嘧啶之间形成两个氢键。}
\end{omfigure}

\section{成环}

\begin{definition}[帕尔–克诺尔合成]\label{def:b3:heterocycles:paal-knorr}
\emph{帕尔–克诺尔合成}\index{帕尔–克诺尔合成}可由 1,4-二羰基化合物制得呋喃
（酸催化），与伯胺或氨反应制得吡咯，或与硫
试剂反应制得噻吩。
\end{definition}

\begin{definition}[汉奇吡啶合成]\label{def:b3:heterocycles:hantzsch}
\emph{汉奇吡啶合成}\index{汉奇吡啶合成}将一个醛、两当量的
$\beta$-酮酯和氨缩合为
1,4-二氢吡啶，该产物可被氧化为吡啶。
\end{definition}

\begin{definition}[费歇尔吲哚合成]\label{def:b3:heterocycles:fischer-indole}
\emph{费歇尔吲哚合成}\index{费歇尔吲哚合成}将醛或酮的
苯腙在酸催化剂作用下加热：腙
互变异构为烯肼，烯肼发生 [3,3] σ 迁移重排；
环化并脱去氨，得到吲哚。
\end{definition}

\begin{omfigure}
\setchemfig{atom sep=1.8em}
\schemestart
\chemfig{*6(-=-=(-NH-[:30]N=[:-30](-[:-90]CH_3)-[:30]CH_3)-=)}
\arrow{->[$\ce{H+}$][互变异构]}[,1.4]
\chemfig{*6(-=-=(-NH-[:30]NH-[:-30](=[:-90]CH_2)-[:30]CH_3)-=)}
\schemestop
\par\bigskip
\schemestart
\arrow{->[{[3,3]}, 然后][环化，$-\ce{NH3}$]}[,2.2]
\chemfig{*6(-=*5(-=(-CH_3)-N(-H)-)-=-=)}
\schemestop

{\small 由丙酮苯腙出发的费歇尔吲哚合成。酸使腙转化为其烯肼互变异构体；[3,3] 重排
断裂 N--N 键，并在与氮相邻的环碳（邻位）与
末端 $\ce{CH2}$ 之间形成 C--C 键；新生成的胺进攻亚胺，
脱去氨后得到
2-甲基吲哚。}
\end{omfigure}

\begin{method}[杂环的逆合成切断]\label{met:b3:heterocycles:disconnect}
\begin{enumerate}
  \item 2,5-二取代吡咯、呋喃或噻吩：切断与杂原子相连的两个键，
        回到 1,4-二酮（帕尔–克诺尔合成）。
  \item 在 C3 和 C5 位带酯基的对称 1,4-二氢吡啶：回到一个醛、
        两当量的 $\beta$-酮酯和氨（汉奇合成）。
  \item 2,3-二取代吲哚：回到苯肼和酮
        $\mathrm{R^2CH_2COR^3}$（费歇尔合成）。
  \item 检查各片段是否容易获得，再考虑
        酮不对称时的区域选择性。
\end{enumerate}
\end{method}

\begin{definition}[生物电子等排体]\label{def:b3:heterocycles:bioisostere}
\emph{生物电子等排体}\index{生物电子等排体}是药物分子中能替代另一基团、同时保持其生物活性的基团，
因为二者在大小、形状及氢键与电荷模式上相近；
例如四氮唑可替代羧酸，
吡啶可替代苯环。
\end{definition}

\begin{inthelab}[帕尔–克诺尔法制备吡咯]
将己烷-2,5-二酮、苯胺与催化量的对甲苯磺酸
在甲苯中于装有迪安–斯塔克分水器的烧瓶里加热回流：
生成的水以共沸物形式被带出，收集在支管中，
从而推动缩合反应进行完全。当不再有水分层析出时，
将溶液用碳酸氢钠水溶液洗涤，
干燥后蒸去溶剂，再将 2,5-二甲基-1-苯基吡咯重结晶。
\end{inthelab}

\begin{safety}{\ghs[0.62cm]{GHS02}\ghs[0.62cm]{GHS06}\ghs[0.62cm]{GHS07}\ghs[0.62cm]{GHS08}\ghs[0.62cm]{GHS09}}
% ledger: ghs4:pyridine, ghs4:phenylhydrazine
吡啶：高度易燃，经各种途径均有害，有强烈气味；
须在通风橱中使用。苯肼：经各种途径均有毒，可能致癌并引起遗传缺陷，
反复接触会损害血液，具有致敏性，对水生生物毒性很大；
须在通风橱中戴手套称量，废液
单独收集。
\end{safety}

\begin{history}[两种经典的成环合成]
埃米尔·费歇尔于 1883 年报道了他的吲哚合成，
同一时期他还研究了糖与嘌呤；阿图尔·汉奇于 1881 年发表了吡啶
合成。两者至今仍在使用：
费歇尔合成用于制备治疗偏头痛的曲普坦类药物，
汉奇合成用于制备治疗高血压的二氢吡啶类钙通道阻滞剂。
\end{history}

\section{练习}

\begin{exercise}[$\star$]\label{exo:b3:heterocycles:1}
计算噁唑、噻唑、嘧啶和嘌呤的 $\pi$ 电子数，并指出
哪些孤电子对属于 $\pi$ 体系。
\end{exercise}

\begin{exercise}[$\star$]\label{exo:b3:heterocycles:2}
按碱性由强到弱排列吡啶、吡咯、咪唑、哌啶和 DMAP，并给出
它们的 $\mathrm pK_a$ 值。
\end{exercise}

\begin{exercise}[$\star$]\label{exo:b3:heterocycles:3}
给出吡咯、呋喃和吲哚单溴代的主要产物。
\end{exercise}

\begin{exercise}[$\star$]\label{exo:b3:heterocycles:4}
给出己烷-2,5-二酮在帕尔–克诺尔反应中分别与单独的酸、
甲胺以及五硫化二磷作用时的产物。
\end{exercise}

\begin{exercise}[$\star\star$]\label{exo:b3:heterocycles:5}
2-氯吡啶在甲醇中于 \qty{65}{\celsius} 与甲醇钠反应，3-氯吡啶
则几乎不反应。解释原因，并给出产物。
\end{exercise}

\begin{exercise}[$\star\star$]\label{exo:b3:heterocycles:6}
给出吡啶的齐齐巴宾反应产物，并画出
其加成中间体。
\end{exercise}

\begin{exercise}[$\star\star$]\label{exo:b3:heterocycles:7}
吡啶 N-氧化物在 C4 位发生硝化。用共振结构解释原因，并说明
如何得到 4-硝基吡啶。
\end{exercise}

\begin{exercise}[$\star\star$]\label{exo:b3:heterocycles:8}
2-羟基吡啶在水中主要以 2-吡啶酮形式存在。画出两种互变异构体，并
解释为何吡啶酮仍具有芳香性。
\end{exercise}

\begin{exercise}[$\star\star$]\label{exo:b3:heterocycles:9}
给出 1,4-二氢吡啶的汉奇吡啶合成所需的组分，该产物
在 C3 和 C5 位带甲酯基，C2 和 C6 位带甲基，
C4 位带 2-硝基苯基。
\end{exercise}

\begin{exercise}[$\star\star\star$]\label{exo:b3:heterocycles:10}
丁酮在费歇尔吲哚合成中与
苯肼反应，得到两种不同的吲哚。写出这两种吲哚，并说明酸的强弱如何改变二者的比例。
\end{exercise}

\begin{exercise}[$\star\star\star$]\label{exo:b3:heterocycles:11}
解释为何 4-氟硝基苯和 2-氟吡啶与亲核试剂
反应比相应的氯代物更快，这与 $\mathrm S_\mathrm N2$ 底物不同。
\end{exercise}

\begin{exercise}[$\star\star\star$]\label{exo:b3:heterocycles:12}
利用本章中的模拟谱图，说明如何通过质子核磁共振谱区分吡啶、
吡咯和呋喃，并预测
噻吩的谱图（化学位移分别为 \qty{7.33}{ppm} 和 \qty{7.12}{ppm}）。
\end{exercise}

\section{问题：构建吲哚}

\begin{problem}[{周末问题——构建吲哚：丙酮苯腙、
烯肼及其 [3,3] 重排、氨的脱去，
以及丁酮所得两种吲哚的选择}]
\label{pb:b3:heterocycles:1}
% ledger: aw:C, aw:H, aw:N, aw:O
题目数据：10.0 克苯肼与过量丙酮反应，并
将所得腙与氯化锌共热。

\medskip\noindent\textbf{第一部分——腙。}
\begin{enumerate}
  \item 写出丙酮苯腙的生成。
  \item 为何属于亚胺型缩合？少量酸为何有利？
  \item 苯肼的哪个氮原子进攻羰基？为什么？
  \item 计算苯肼的摩尔质量及所用物质的量。
  \item 计算腙的理论质量。
\end{enumerate}

\medskip\noindent\textbf{第二部分——烯肼与 [3,3] 步骤。}
\begin{enumerate}[resume]
  \item 画出烯肼互变异构体。
  \item 构成 [3,3] 重排环状过渡态的是哪六个原子？
  \item 哪一个键断裂，哪一个键形成？
  \item 参与反应的电子数为多少？该步骤是否为热允许的？
  \item 氯化锌起什么作用？
  \item 为何之后必须恢复苯环的芳香性？
\end{enumerate}

\medskip\noindent\textbf{第三部分——成环。}
\begin{enumerate}[resume]
  \item 哪个氮原子进攻哪个碳原子，从而闭合五元环？
  \item 离去的小分子是什么？最后一步的驱动力是什么？
  \item 写出由腙生成 2-甲基吲哚的总反应方程式。
  \item 计算 2-甲基吲哚的摩尔质量。
  \item 计算 2-甲基吲哚的理论质量。
  \item 2-甲基吲哚在哪个位置发生亲电取代？为什么？
\end{enumerate}

\medskip\noindent\textbf{第四部分——丁酮。}
\begin{enumerate}[resume]
  \item 画出丁酮苯腙的两种烯肼。
  \item 写出它们分别生成的两种吲哚。
  \item 哪种烯肼取代程度更高？哪一种在强酸条件下生成？
  \item 弱酸条件下哪种吲哚占优势？为什么？
  \item 为何乙醛苯腙不能很好地生成吲哚本身？
  \item 区域选择性问题对药物合成意味着什么？
  \item 给出结果：由 10.0 克
        苯肼得到的 2-甲基吲哚的理论质量。
\end{enumerate}
\end{problem}
